% 第九卷 5.3 节：以PPO为优化器的RLHF三阶段流程（跨栏，约160 mm宽）。
% 上排：第一阶段有监督微调、第二阶段奖励模型训练；下排：第三阶段PPO。
% 蓝色：数据；黄色实线框：参与训练的模型；灰色虚线框：冻结的模型；红色：奖励与优势等训练信号。
\begin{tikzpicture}[
  x=1mm,y=1mm,
  figure picture,
  data/.style={figure node,draw=FigureInput,fill=FigureInputFill,minimum height=8mm,
    minimum width=18mm,inner xsep=2pt},
  train/.style={figure node,draw=FigureModel,fill=FigureModelFill,minimum height=8mm,
    minimum width=20mm,inner xsep=2pt,line width=.9pt},
  frozen/.style={figure node,draw=FigureMuted,fill=white,minimum height=8mm,
    minimum width=20mm,inner xsep=2pt,dash pattern=on 2pt off 1.5pt},
  signal/.style={figure node,draw=FigureOutput,fill=FigureOutputFill,minimum height=8mm,
    minimum width=18mm,inner xsep=2pt},
  stage/.style={figure subgroup},
  stagehead/.style={figure heading,anchor=north west}
]
  % ---------- 阶段一：有监督微调 ----------
  \draw[stage] (0,44) rectangle (46,76);
  \node[stagehead] at (1.5,75) {阶段一：有监督微调};
  \node[data] (demo) at (12,56) {人工示范};
  \node[train] (sft) at (34,56) {SFT模型};
  \draw[figure flow] (demo.east) -- (sft.west);

  % ---------- 阶段二：奖励模型训练 ----------
  \draw[stage] (52,44) rectangle (160,76);
  \node[stagehead] at (53.5,75) {阶段二：奖励模型训练};
  \node[frozen] (samp) at (65,56) {SFT采样\\$K$个回答};
  \node[data] (rank) at (91,56) {标注者排序};
  \node[data] (pairs) at (117,56) {展开为\\$\binom K2$对};
  \node[train] (rm) at (146,56) {奖励模型$r_{\bm\phi}$};
  \draw[figure flow] (samp.east) -- (rank.west);
  \draw[figure flow] (rank.east) -- (pairs.west);
  \draw[figure flow] (pairs.east) -- (rm.west);
  \draw[figure flow,draw=FigureMuted] (sft.east) -- (samp.west);

  % ---------- 阶段三：PPO ----------
  \draw[stage] (0,0) rectangle (160,38);
  \node[stagehead] at (1.5,37) {阶段三：PPO强化学习};
  \node[data,minimum width=12mm] (x) at (9,18) {提示$x$};
  \node[train] (pol) at (32,18) {策略$\pi_{\bm\theta}$\\（初值SFT）};
  \node[data,minimum width=13mm] (y) at (54,18) {回答$y$};
  \node[frozen] (rmf) at (78,26) {奖励模型$r_{\bm\phi}$};
  \node[frozen] (ref) at (78,10) {参考$\pi_{\mathrm{ref}}$\\（SFT）};
  \node[signal] (rt) at (106,18) {逐词元\\奖励$r_t$};
  \node[train] (val) at (129,18) {价值$V_{\bm\omega}$\\（初值RM）};
  \node[signal,minimum width=14mm] (ppo) at (151,18) {GAE与\\PPO};
  \draw[figure flow] (x.east) -- (pol.west);
  \draw[figure flow] (pol.east) -- (y.west);
  \draw[figure flow] (y.east) -- ++(4,0) |- (rmf.west);
  \draw[figure flow] (y.east) -- ++(4,0) |- (ref.west);
  \draw[figure flow] (rmf.east) to[out=0,in=180] ([yshift=2mm]rt.west);
  \draw[figure flow] (ref.east) to[out=0,in=180] ([yshift=-2mm]rt.west);
  \draw[figure flow] (rt.east) -- (val.west);
  \draw[figure flow] (val.east) -- (ppo.west);
  \draw[figure flow,draw=FigureOutput] (ppo.south) -- (151,3) -| node[figure note,pos=0.17,above] {更新$\bm\theta$与$\bm\omega$} (pol.south);
\end{tikzpicture}
